% ===========================================================================
\chapter{Finding things}\label{ch:search}

In a project of four files, the model can read everything. In a project of
four hundred, it can't, and shouldn't try: every file it reads fills up
the conversation (and the bill). A programmer in a big project doesn't
read it all either. They search: where is \code{days_late} defined, and
who calls it? That is our next tool.

The search takes a \emph{regular expression}, a small language for
patterns of text: \code{days_late} matches those letters anywhere in a
line, \code{fine\textbackslash(} matches \code{fine(}, and
\code{\textasciicircum class} matches a line that starts with
\code{class}. Models know regular expressions well, which makes them a
good choice of input: a tool is easiest to use when its input is
something the model has seen a million times. Nex's \code{text/Regex}
library supplies them. Add it to the files \code{tools.nex} interns:

\begin{tightaddbefore}{tools.nex}{before intern data/Json}{intern data/Json}
intern text/Regex
\end{tightaddbefore}

Each match comes back as \code{file:line: text}, the format of
\code{grep}, which the model knows too, and which tells it where to look
next with \code{read_file}.

\begin{append}{tools.nex}
class Search inherit Tool
create
  make(ws: Workspace) do
    super.make(ws)
  end

feature
  MOST_MATCHES = 100

  name(): String do
    result := "search"
  end

  description(): String do
    result := "Search the text files in the workspace for lines matching a regular "
      + "expression (Java syntax). Each match comes as `file:line: text`. Use it to "
      + "find where something is defined or used without reading every file."
  end

  parameters(): Map[String, Any] do
    result := {"pattern": parameter("string", "The regular expression, e.g. `days_late` or `fine\\(`."),
      "path": parameter("string", "The folder to search, relative to the workspace. Default: `.`.")}
  end

  required(): Array[String] do
    result := ["pattern"]
  end

  describe(input: Json_Value): String do
    result := "search " + string_argument(input, "pattern", "?")
    if string_argument(input, "path", ".") /= "." then
      result := result + " in " + string_argument(input, "path", ".")
    end
  end
\end{append}

(In \code{"fine\textbackslash\textbackslash("}, the two backslashes are one
backslash in the string: this is an ordinary string, not a raw one.)

A pattern can be wrong: \code{fine(} has an opening bracket with no
closing one. \code{Regex.compile} then raises an exception, which we turn
into an error the model can read. The \code{regex} variable is a
\code{?Regex}, \code{nil} unless the compiling worked:

\begin{continued}{tools.nex}
  run(input: Json_Value): Tool_Outcome do
    let pattern: String := string_argument(input, "pattern", "")
    let path: String := string_argument(input, "path", ".")
    let problem: String := path_problem(path, "folder")
    if problem = "" and pattern = "" then
      problem := "the pattern is empty"
    end
    let regex: ?Regex := nil
    if problem = "" then
      do
        regex := create Regex.compile(pattern)
      rescue
        problem := "not a valid regular expression: " + exception
      end
    end
    if ?regex as r then
      result := create Tool_Outcome.ok(matches(r, workspace.files(path)))
    else
      result := create Tool_Outcome.error(problem)
    end
  end
\end{continued}

The search goes through every file, line by line. Not every file is text:
reading a picture as text fails, and such a file is simply skipped. A
common word could match thousands of lines, so after 100 matches the rest
are only counted. Telling the model how many it didn't see lets it narrow
the search, instead of believing it has seen them all.

\begin{continued}{tools.nex}
  private feature
    matches(regex: Regex, files: Array[String]): String do
      let found: Array[String] := []
      let more: Integer := 0
      across files as file do
        let text: String := ""
        do
          text := workspace.read(file)
        rescue
          text := ""     -- not text (a picture, say): nothing to find in it
        end
        let number: Integer := 1
        across lines_of(text) as line do
          if regex.find(line) /= nil then
            if found.length < MOST_MATCHES then
              found.add(file + ":" + number.to_string + ": " + line.trim())
            else
              more := more + 1
            end
          end
          number := number + 1
        end
      end
      if found.is_empty then
        result := "(no matches)"
      else
        result := shortened(joined_lines(found), output_limit())
        if more > 0 then
          result := result + "... and " + count_text(more, "more match") + "\n"
        end
      end
    end
end
\end{continued}

\code{regex.find(line)} gives the part of the line that matches, or
\code{nil} if none does.

\begin{tryit}{try_search.nex}
intern Tools

let ws: Workspace := fresh_workspace(user_path("sample_project"), user_path("scratch"))
let search: Search := create Search.make(ws)
across ['{"pattern": "days_late"}', '{"pattern": "^class"}', '{"pattern": "fine\\("}',
    '{"pattern": "fine("}', '{"pattern": "overdue"}'] as input do
  let arguments: Json_Value := json_value(input)
  let outcome: Tool_Outcome := search.run(arguments)
  console.print_line("> " + search.describe(arguments))
  console.print_line((when outcome.is_error then "error: " else "" end) + outcome.text)
end
\end{tryit}
\tryoutput{try_search}

The third pattern shows the cost of escaping in layers: the regular
expression \code{fine\textbackslash(} is written \code{fine\textbackslash\textbackslash(}
inside JSON. The model writes JSON for a living and gets this right; we
had to think about it.
